\chapter{本科量子散射理论}

\begin{prereqbox}
定态 Schrödinger 方程、球坐标分离变量、概率流和 Fourier 变换。本章只讨论非相对论势散射，不使用场论语言。
\end{prereqbox}

\section{从实验计数到散射截面}
设短程势局域在原点附近，入射粒子动量 $\hbar k\hat{\vb z}$。远离散射中心时
\[
\psi(\vb r)\sim e^{ikz}+f(\theta,\phi)\frac{e^{ikr}}r.
\]
第一项是入射平面波，第二项是出射球面波。概率流
\[
\vb j=\frac{\hbar}{2mi}(\psi^*\nabla\psi-\psi\nabla\psi^*)
\]
对入射波给
\[
\vb j_{\rm in}=\frac{\hbar k}{m}\hat{\vb z}.
\]
对远区散射波 $\psi_s=f e^{ikr}/r$，保留 $1/r^2$ 的径向通量：
\[
j_{s,r}=\frac{\hbar k}{m}\frac{|f(\theta,\phi)|^2}{r^2}.
\]
穿过面积元 $r^2\dd\Omega$ 的散射粒子数率除以入射流强，定义微分截面：
\[
\boxed{\dv{\sigma}{\Omega}
=\frac{j_{s,r}r^2}{j_{\rm in}}=|f(\theta,\phi)|^2}.
\]
$f$ 的量纲是长度，截面的量纲是面积。总截面为
\[
\sigma=\int |f|^2\dd\Omega.
\]
若末态波数不同，还需乘速度比 $v_f/v_i$；本章主要讨论弹性散射。

\section{自由 Green 函数与出射边界条件}
定态方程写成
\[
(E-H_0)\ket\psi=V\ket\psi,\qquad
H_0=-\frac{\hbar^2}{2m}\nabla^2.
\]
形式逆算符在自由能谱上奇异，因此定义
\[
G_0^{(+)}(E)=\frac1{E-H_0+i0}.
\]
坐标表象满足
\[
\left(\nabla^2+k^2\right)G_k^{(+)}(\vb r-\vb r')
=\frac{2m}{\hbar^2}\delta^{(3)}(\vb r-\vb r').
\]
先求 Helmholtz 方程的基本解。对 $R\ne0$，球对称解为
\[
g(R)=\frac{Ae^{ikR}+Be^{-ikR}}R.
\]
$+i0$ 选择向外传播的 $e^{ikR}$。为定系数，对半径 $\epsilon$ 小球积分：
\[
\int_{R<\epsilon}\nabla^2g\,\dd^3R
=4\pi\epsilon^2\pdv{g}{R}\bigg|_\epsilon
\longrightarrow-4\pi A.
\]
与右端 $2m/\hbar^2$ 比较，得到
\[
\boxed{
G_k^{(+)}(\vb r,\vb r')
=-\frac{m}{2\pi\hbar^2}
\frac{e^{ik|\vb r-\vb r'|}}{|\vb r-\vb r'|}}.
\]
于是
\[
\boxed{\psi^{(+)}(\vb r)
=e^{i\vb k\cdot\vb r}
+\int G_k^{(+)}(\vb r,\vb r')V(\vb r')
\psi^{(+)}(\vb r')\dd^3r'}
\]
即 Lippmann--Schwinger 方程。它不是近似，而是带出射边界条件的 Schrödinger 方程积分形式。

\section{由远区展开读出散射振幅}
当 $r\gg r'$，
\[
|\vb r-\vb r'|
=r-\hat{\vb r}\cdot\vb r'+O(r'^2/r),
\]
\[
\frac{e^{ik|\vb r-\vb r'|}}{|\vb r-\vb r'|}
\simeq\frac{e^{ikr}}r e^{-ik\hat{\vb r}\cdot\vb r'}.
\]
令 $\vb k'=k\hat{\vb r}$，与远区定义比较：
\[
\boxed{
f(\vb k',\vb k)
=-\frac{m}{2\pi\hbar^2}
\int e^{-i\vb k'\cdot\vb r'}
V(\vb r')\psi^{(+)}_{\vb k}(\vb r')\dd^3r'}.
\]
精确波函数仍出现在右端，所以这还是隐式表达式。

\section{Born 级数与一阶 Born 近似}
把积分方程反复代回：
\[
\ket{\psi^{(+)}}=
\ket{\vb k}+G_0^{(+)}V\ket{\vb k}
+G_0^{(+)}VG_0^{(+)}V\ket{\vb k}+\cdots.
\]
若散射势足够弱或粒子能量足够高，第一轮中取
$\psi^{(+)}\simeq e^{i\vb k\cdot\vb r}$：
\[
\boxed{
f_B(\vb q)=-\frac{m}{2\pi\hbar^2}
\int e^{-i\vb q\cdot\vb r}V(\vb r)\dd^3r},
\qquad
\vb q=\vb k'-\vb k.
\]
弹性散射中
\[
q=2k\sin\frac\theta2.
\]
因此角分布直接测量势能的空间 Fourier 分量。粗略判据是波函数在势区内不被明显改造；对尺度 $a$、强度 $V_0$ 的势，常要求
$m|V_0|a^2/\hbar^2\ll1$，高能时条件可放宽。

\begin{examplebox}[title=Gaussian 势的 Born 散射]
取 $V(r)=V_0e^{-r^2/a^2}$。三维 Fourier 积分分离：
\[
\int e^{-i\vb q\cdot\vb r}e^{-r^2/a^2}\dd^3r
=\prod_{\alpha=x,y,z}\int_{-\infty}^{\infty}
e^{-r_\alpha^2/a^2-iq_\alpha r_\alpha}\dd r_\alpha.
\]
完成平方后每个积分为
$\sqrt\pi a\,e^{-a^2q_\alpha^2/4}$，故
\[
f_B(q)=-\frac{mV_0\sqrt\pi a^3}{2\hbar^2}
e^{-a^2q^2/4},
\]
\[
\dv\sigma\Omega
=\frac{m^2V_0^2\pi a^6}{4\hbar^4}
e^{-a^2q^2/2}.
\]
$ka\gg1$ 时大角度对应大 $q$，被指数压低，散射集中在前向。
\end{examplebox}

\begin{examplebox}[title=Yukawa 势与 Rutherford 极限]
取
\[
V(r)=\alpha\frac{e^{-\kappa r}}r.
\]
选 $\vb q$ 为极轴，
\[
\int e^{-i\vb q\cdot\vb r}V(r)\dd^3r
=2\pi\alpha\int_0^\infty r e^{-\kappa r}\dd r
\int_{-1}^{1}e^{-iqru}\dd u.
\]
角积分为 $2\sin(qr)/(qr)$，所以
\[
=\frac{4\pi\alpha}{q}
\int_0^\infty e^{-\kappa r}\sin(qr)\dd r
=\frac{4\pi\alpha}{q^2+\kappa^2}.
\]
于是
\[
f_B(q)=-\frac{2m\alpha}{\hbar^2(q^2+\kappa^2)}.
\]
令 $\kappa\to0$ 并用 $q=2k\sin(\theta/2)$：
\[
\dv\sigma\Omega
=\frac{m^2\alpha^2}{4\hbar^4k^4\sin^4(\theta/2)}.
\]
这给 Rutherford 角依赖，但纯 Coulomb 势是长程势，普通出射渐近式会出现对数相位；严格处理不能仅靠“令屏蔽消失”替代。
\end{examplebox}

\section{中心势的分波展开}
轴对称入射允许展开
\[
\psi(r,\theta)=\sum_{l=0}^\infty A_lR_l(r)P_l(\cos\theta).
\]
令 $u_l=rR_l$，径向方程为
\[
\frac{\dd^2u_l}{\dd r^2}
+\left[k^2-\frac{2m}{\hbar^2}V(r)
-\frac{l(l+1)}{r^2}\right]u_l=0.
\]
离心势在低能时阻止高 $l$ 波进入势区。

在势外，解是球 Bessel 函数的组合：
\[
R_l(r)=C_l[\cos\delta_l\,j_l(kr)-\sin\delta_l\,n_l(kr)].
\]
利用远区渐近式，
\[
R_l(r)\sim\frac{C_l}{kr}
\sin\left(kr-\frac{l\pi}{2}+\delta_l\right).
\]
$\delta_l$ 就是势能相对于自由波引入的相移。平面波展开为
\[
e^{ikz}=\sum_l i^l(2l+1)j_l(kr)P_l(\cos\theta).
\]
匹配每个分波的入射部分后得到
\[
\boxed{
f(\theta)=\frac1k\sum_{l=0}^\infty
(2l+1)e^{i\delta_l}\sin\delta_l\,P_l(\cos\theta)}.
\]
利用 Legendre 多项式正交性
\[
\int P_lP_{l'}\dd\Omega=\frac{4\pi}{2l+1}\delta_{ll'},
\]
总截面为
\[
\boxed{\sigma=\frac{4\pi}{k^2}
\sum_l(2l+1)\sin^2\delta_l}.
\]
每个分波的贡献受上限
$\sigma_l\le4\pi(2l+1)/k^2$ 约束，这是幺正性限制。

\section{光学定理}
前向 $P_l(1)=1$，所以
\[
\operatorname{Im}f(0)
=\frac1k\sum_l(2l+1)\sin^2\delta_l.
\]
与总截面对比：
\[
\boxed{\sigma_{\rm tot}=\frac{4\pi}{k}\operatorname{Im}f(0)}.
\]
它看似把“所有角度的损失”交给前向振幅，实质是入射波与前向散射波的干涉恰好扣除离开原束流的概率，因而是概率流守恒的整体表达。

\section{低能 $s$ 波、散射长度与浅束缚态}
若势程为 $R$ 且 $kR\ll1$，离心势使 $l>0$ 的相移按更高次幂消失，只剩
\[
f_0(k)=\frac1{k\cot\delta_0-ik}.
\]
有效程展开为
\[
k\cot\delta_0=-\frac1{a_s}+\frac12r_ek^2+\cdots,
\]
故零能极限
\[
f_0\to-a_s,\qquad \sigma\to4\pi a_s^2.
\]
若 $a_s$ 很大且为正，$f_0$ 在复 $k$ 平面靠近 $k=i/a_s$ 有极点，对应能量
\[
E_b\simeq-\frac{\hbar^2}{2ma_s^2}
\]
的浅束缚态。调节势使束缚态刚好穿过阈值时，$a_s$ 发散并出现低能共振。

\begin{examplebox}[title=硬球低能散射]
硬球半径为 $a$，边界条件 $u_l(a)=0$。低能只保留 $l=0$：
\[
u_0(r)\propto\sin(kr+\delta_0).
\]
故
\[
\sin(ka+\delta_0)=0,\qquad \delta_0=-ka
\]
（选择连续趋于零的分支）。于是
\[
a_s=-\lim_{k\to0}\frac{\tan\delta_0}{k}=a,
\qquad
\sigma\to4\pi a^2.
\]
它是经典几何截面 $\pi a^2$ 的四倍，来源于低能波的相干衍射。
\end{examplebox}

\begin{examplebox}[title=球形方阱的散射长度]
取 $V=-V_0$（$r<R$），外部为零。零能 $s$ 波在阱内
\[
u_{\rm in}=A\sin(Kr),\qquad
K=\sqrt{2mV_0}/\hbar,
\]
外部零能径向方程给
\[
u_{\rm out}=B(r-a_s).
\]
在 $R$ 处匹配对数导数：
\[
K\cot(KR)=\frac1{R-a_s}.
\]
解得
\[
\boxed{a_s=R-\frac{\tan(KR)}K}.
\]
当 $KR=(n+\tfrac12)\pi$，散射长度发散，表示新的 $s$ 波束缚态恰好到达阈值。
\end{examplebox}

\section{相同粒子的散射}
不可分辨粒子有两条无法区分的路径：散到 $\theta$，或交换标签后散到
$\pi-\theta$。因此必须先叠加振幅，再取模平方：
\[
f_\pm(\theta)=f(\theta)\pm f(\pi-\theta).
\]
无自旋玻色子取加号；相同且自旋极化的费米子取减号。不可分辨的末态只应计数一次，故可积半球，也可在全立体角积分后乘 $1/2$：
\[
\sigma_\pm
=\frac12\int_{4\pi}|f_\pm(\theta)|^2\dd\Omega.
\]
两个未极化自旋 $1/2$ 费米子的一个单态通道配空间对称振幅，三个三重态通道配空间反对称振幅：
\[
\dv{\sigma}{\Omega}
=\frac14|f_+|^2+\frac34|f_-|^2.
\]
低能 $s$ 波近似中 $f(\theta)$ 与角度无关，故
$f_-=0$：自旋极化费米子的 $s$ 波散射被 Pauli 原理抑制。交换对称性、单态/三重态和重复计数的完整推导见第
\ref{chap:identical-particles} 章第
\ref{sec:identical-scattering-counting} 节。

\section{与凝聚态和实验的连接}
杂质势的一阶 Born 振幅含 $V(\vb q)$；黄金规则给
\[
\frac1\tau\propto\sum_{\vb k'}|V(\vb k'-\vb k)|^2
\delta(E_{\vb k'}-E_{\vb k}).
\]
输运寿命还多出 $1-\cos\theta$，因为前向散射几乎不松弛电流。中子和 X 射线散射中，动量转移依旧是 $\vb q=\vb k'-\vb k$，测得的强度则由相应密度或自旋关联函数控制。

\begin{interviewbox}
\textbf{两分钟回答散射理论：}先写远区波函数并由概率流得到
$\dd\sigma/\dd\Omega=|f|^2$。一般势可把 Schrödinger 方程写成带出射 Green 函数的 Lippmann--Schwinger 方程；弱势时第一次迭代给 Born 近似，振幅是势的 Fourier 变换。中心势更适合分波法，每个角动量通道只引入相移 $\delta_l$；低能因离心势只剩 $s$ 波，散射由一个长度 $a_s$ 控制。光学定理则是总概率守恒。
\end{interviewbox}

\begin{selftestbox}
\begin{enumerate}
\item 从球面波概率流检查 $f$ 的量纲。
\item $+i0$ 为什么选择出射而不是入射边界条件？
\item 推导 $q=2k\sin(\theta/2)$。
\item Born 近似为何在强吸引势接近束缚态阈值时失败？
\item 从分波振幅推出总截面。
\item 光学定理如何体现入射与前向散射波的干涉？
\item 为什么低能自旋极化费米子没有 $s$ 波散射？
\item 输运寿命为什么含 $1-\cos\theta$？
\end{enumerate}
\end{selftestbox}
